% Chapter 4 -- KV cache and the serving seam.
% Sources actually read for this chapter:
%   llm-serving/inferenceLM/engine/{lm_engine.py,transformer_lm_adapter.py}
%   llm-serving/tests/{transformer_lm_adapter_test.py,trained_model_integration_test.py}
%   cs336-basics/cs336_basics/lm.py, transfromer/{transformer,multiheads_attention,
%     scaled_dot_prod_attention}.py
%   tests/test_kv_cache.py, cs336_systems/experiments/{kv_cache_gate,bench_serving}.py
%   artifacts/serving_benchmark{,_kv}.json, artifacts/NIGHT_REPORT.md,
%   run_logs/post_train_window.log
% Cross-repo files are quoted by hand (never \lstinputlisting): the handout must build
% from a standalone clone of LLM-system-project, where ../../../llm-serving/ does not exist.

\section{KV cache 与 serving 集成：把两个仓库接起来}
\label{sec:b-kv}

到这一步为止，两个仓库互不相识。\code{LLM-system-project} 训出了一个 0.19B 的
\code{TransformerLM}，它的公开接口就一个：\code{forward(x) -> logits}，
输入 \code{(batch, seq)} 的 token id，输出 \code{(batch, seq, vocab)} 的 logits。
\code{llm-serving} 那边的 \code{LMEngine} 是照着 HuggingFace 的形状写的：
它调 \codeb{model(input_ids=..., past_key_values=..., use_cache=True)}，
从返回值上取 \code{.logits} 和 \code{.past_key_values}，还会去读 \code{model.config}。
两边一样都对不上。

摆在面前的有两条路：改 \code{LMEngine} 让它容纳两种模型，或者造一个 adapter 把
\code{TransformerLM} 包成 HF 的形状。这一章的第一件事，是说明\cnemph{"只补四样东西"这个约束
是怎么把设计钉死的}——一旦确认 \code{LMEngine} 依赖的 HF 约定\cnemph{恰好只有四样}，
adapter 就不再是"一个兼容层"这种含糊的东西，而是一个有明确边界、可以被逐条检查的对象。

第二件事是 KV cache 本身：它怎么做成一个对训练路径零风险的纯增量改动，怎么被测，
以及\cnemph{它实测的收益和我预期的完全不一样}（\secref{sec:b-kv-flat}，本章的重点）。

\subsection{\codet{LMEngine} 到底依赖 HF 的什么}
\label{sec:b-kv-four}

\code{LMEngine} 是故意手动驱动模型的，不用 \code{.generate()}——这个决定和它的理由属于姊妹篇，
见 \extref{a:why-not-generate}。副作用是它把 HF 的调用约定摊开写在了自己的代码里，
于是可以被逐行数出来。数的结果是\cnemph{四样}，一样不多：

\begin{table}[htbp]
\centering
\small
\setlength{\tabcolsep}{4pt}
\begin{tabular}{p{0.32\textwidth}p{0.09\textwidth}p{0.48\textwidth}}
\toprule
\codet{LMEngine} 要的 HF 约定 & 本项目 有无 & adapter 怎么补 \\
\midrule
\codet{model.eval()} & 有 &
  不用补。\codet{TransformerLMAdapter} 继承 \codet{nn.Module}，
  \codet{.eval()} 会递归到被包的模型。 \\
\addlinespace
\codeb{model(input_ids, past_key_values, use_cache)} 返回带 \codet{.logits} 的对象 & 部分有 &
  \codet{forward} 接受这三个关键字参数（\codet{use\_cache} 收下但不使用），
  把结果装进 \codet{CausalLMOutput} 这个 \codet{NamedTuple}。 \\
\addlinespace
返回值还要有 \codet{.past\_key\_values}，且能原样喂回下一步 & 没有 &
  同一个 \codet{NamedTuple} 的第二个字段。默认模式下它装的是 token 前缀，
  \codet{use\_kv\_cache=True} 时装的是逐层 $(K,V)$。 \\
\addlinespace
\codet{model.config.n\_positions} & 没有 &
  \codet{AdapterConfig} 这个 dataclass 的字段，构造时由 \codet{context\_length} 参数填。 \\
\addlinespace
\codeb{model.config.eos_token_id} & 没有 &
  同一个 dataclass 的字段，默认 50256（GPT-2 的 \codet{<|endoftext|>}）。 \\
\bottomrule
\end{tabular}
\caption{\codet{LMEngine} 依赖的四样 HF 约定（第二和第三行是同一次调用的两个返回字段，
所以数成四样）$\times$ \codet{TransformerLM} 有没有 $\times$ adapter 怎么补。
出处：逐行读 \codeb{llm-serving/inferenceLM/engine/lm_engine.py}，
四个依赖点分别在第 12 行（\codet{eval()}）、第 36 和 60 行（模型调用）、
第 38 和 40 行（\codet{.logits} / \codet{.past\_key\_values}）、
第 89 行（\codet{n\_positions}）、第 119 行（\codet{eos\_token\_id}）。}
\label{tab:four-conventions}
\end{table}

这个"只有四样"的清单，以及为什么是这四样而不是别的，在姊妹篇里有完整的推导：
\extref{a:lmengine-four-methods}。这一章只用它的结论。

\begin{checkpoint}{four_hf_conventions}{为什么"数清楚是四样"比"写一个兼容层"重要}
不看代码回答：
\begin{enumerate}[label=(\alph*)]
  \item 如果不先数清楚，最可能写出什么样的 adapter？
  \item \code{use_cache} 这个参数 adapter 收下了却不用，为什么不干脆删掉？
  \item 少补了其中一样，会在什么时候暴露？
\end{enumerate}
\hint 想想 \code{config} 这个对象在 HF 里有多少个字段。

\ifshowanswers
\tcblower
\answer (a) 大概率会写成"尽量像 HF"——把 \code{PretrainedConfig} 的一堆字段照抄一遍，
返回值也做成一个字段齐全的 dataclass。那样 adapter 的正确性就没有边界了：
你没法回答"它够不够"，只能回答"它像不像"。数清楚之后，adapter 的规格变成一个五行的清单，
每一行都能对应到 \code{lm_engine.py} 的一个具体行号，也就能被逐条测。
\code{AdapterConfig} 的 docstring 里那句 "the subset of a HuggingFace \code{PretrainedConfig}
that \code{LMEngine} actually reads" 就是这个意思。

(b) 因为它是\cnemph{调用方的}约定，不是被调方的选择。\code{LMEngine.decode()} 里写死了
\code{use_cache=True}；adapter 如果不接这个关键字，那一行就是 \code{TypeError}。
收下但不使用，等于把"这个开关在我们这边没有意义"这件事写进签名里，
而不是逼调用方为我们改代码。

(c) 在\cnemph{请求时}，不是构造时。四样里三样是属性访问（\code{config.n_positions}、
\code{config.eos_token_id}、\code{outputs.past_key_values}），Python 不会在
\codeb{TransformerLMAdapter(model, ...)} 这一步检查它们，只会在第一次真的处理请求时抛
\code{AttributeError}。这就是为什么 \codeb{transformer_lm_adapter_test.py} 里专门有一个
只断言两个 config 字段的测试——它的 docstring 明说了这一点。

\whereincode \codeat{llm-serving/inferenceLM/engine/lm\_engine.py}{89} 与 \codeat{llm-serving/inferenceLM/engine/lm\_engine.py}{119}；
测试在 \codeat{llm-serving/tests/transformer\_lm\_adapter\_test.py}{86}。
\fi
\end{checkpoint}

\subsection{三个类，和一条不 import 的规矩}
\label{sec:b-kv-adapter}

整个 adapter 是三个类，连同 39 行模块 docstring 一共 151 行。

\begin{lstlisting}[style=hlplain]
@dataclass
class AdapterConfig:
    """The subset of a HuggingFace `PretrainedConfig` that `LMEngine` actually reads."""
    n_positions: int
    eos_token_id: int
    vocab_size: Optional[int] = None


class CausalLMOutput(NamedTuple):
    logits: torch.Tensor
    past_key_values: torch.Tensor
\end{lstlisting}
\noindent 来源：\codeb{llm-serving/inferenceLM/engine/transformer_lm_adapter.py}:50--67，
手抄（跨仓库文件，不能用 \codeb{\textbackslash lstinputlisting}）。\code{CausalLMOutput}
的 docstring 已省略，正文在 \secref{sec:b-kv-recompute} 引用。

第三个类 \code{TransformerLMAdapter} 继承 \code{nn.Module}，构造时收下一个已经建好的模型、
一个 \code{context_length}、一个 \code{eos_token_id}，以及一个 \code{use_kv_cache} 开关。
它的 \code{forward} 就是表~\ref{tab:four-conventions} 第二行那个签名。

真正值得停一下的是它\cnemph{没有}做的一件事：\code{transformer_lm_adapter.py}
从头到尾没有 \code{import cs336_basics}。它的类型标注是 \code{model: nn.Module}，
docstring 写着 "Typed loosely on purpose"。这有三个后果，三个都是当初想要的：

\begin{enumerate}[label=(\arabic*), leftmargin=1.8em]
  \item \cnemph{\codet{inferenceLM} 对训练项目零依赖。} 装了 \code{llm-serving} 的人
        不需要装 \code{cs336_basics}；两个仓库仍然可以各自独立 clone、独立测试。
        唯一需要两边都在的是 \codeb{tests/trained_model_integration_test.py}，
        它用 \code{pytest.importorskip} 在缺一半的环境里自己跳过
        （\codeat{llm-serving/tests/trained\_model\_integration\_test.py}{37}）。
  \item \cnemph{adapter 可以对着 stub 测。} 只要是 \code{(batch, seq) -> logits} 的
        \code{nn.Module} 都能塞进去——\secref{sec:b-kv-stub} 会看到塞进去的其实是 GPT-2。
  \item \cnemph{绕开了 \codet{llm-serving/CLAUDE.md} 的 Forbidden 约束。} 那份文件写着
        "不许在 \code{inferenceLM/} 和 \code{tests/} 之外改文件"、
        "不许不加论证地引入新依赖"。一个 import 训练项目的 adapter 会同时违反两条；
        一个只收 \code{nn.Module} 的 adapter 一条都不违反。
\end{enumerate}

\begin{checkpoint}{adapter_not_branch}{为什么造 adapter，不是在 \codet{LMEngine} 里加分支}
不看代码回答：
\begin{enumerate}[label=(\alph*)]
  \item 在 \code{LMEngine} 里加一个 \codeb{if isinstance(model, TransformerLM)} 分支，
        会带来什么具体后果？
  \item adapter 这个方案的代价是什么？（它不是免费的）
\end{enumerate}

\ifshowanswers
\tcblower
\answer (a) 三个后果。第一，\code{LMEngine} 要 import 训练项目，两个仓库从此绑死，
\code{llm-serving} 的 \code{uv sync} 得把 \code{einops}/\code{tiktoken} 一起拉进来。
第二，\code{prefill()}/\code{decode()}/\code{inference()} 三个方法里每一处碰
\code{past_key_values} 或 \code{config} 的地方都要开分支，而这三个方法正是将来要改成
continuous batching 的地方——分支会被复制到每一个新特性里。第三，
"\code{LMEngine} 支持哪些模型"这个问题从此没有静态答案，只能靠读分支。

(b) 代价是多了一层间接，以及\cnemph{一个字段名装了两种完全不同的东西}——
\code{past_key_values} 在默认模式下是 token 前缀，在 \code{use_kv_cache=True} 下是逐层 $(K,V)$。
这是真实的可读性代价，\secref{sec:b-kv-recompute} 的例子把两种模式并排放，就是为了让这一点无处可藏。
它之所以可以接受，是因为 \code{LMEngine} 把这个字段当\cnemph{不透明}对象处理：
它只负责把上一步的返回值原样传回去，从不解释里面装的是什么。
\fi
\end{checkpoint}

\begin{checkpoint}{loose_typing_on_purpose}{松类型标注是偷懒还是设计}
不看代码回答：
\begin{enumerate}[label=(\alph*)]
  \item \code{model: nn.Module} 而不是 \code{model: TransformerLM}，放弃了什么？换来了什么？
  \item 有没有一个地方，adapter 还是做了运行时检查？
\end{enumerate}

\ifshowanswers
\tcblower
\answer (a) 放弃的是静态类型检查——\code{mypy} 不会告诉你塞进去的东西有没有正确的 forward 签名，
签名不对要到第一次 forward 才炸。换来的是 \secref{sec:b-kv-adapter} 里那三条：零依赖、
可对 stub 测、不违反仓库约束。这是一次明确的取舍，不是忘了写标注：docstring 里
"Typed loosely on purpose" 那半句就是在说这件事。

(b) 有，而且只有一处：\code{use_kv_cache=True} 时会 \codeb{hasattr(model, "forward_with_cache")}，
不满足就在\cnemph{构造时} \code{raise ValueError}，错误信息里直接告诉你正确的签名，
以及"留 False 就能用任何 \code{(batch, seq) -> logits} 模块"。
选择在这里做检查而不是别处，是因为这是唯一一个"松类型会导致运行到一半才失败"的地方——
其余的能力缺失在第一次 forward 就会暴露。

\whereincode \codeat{llm-serving/inferenceLM/engine/transformer\_lm\_adapter.py}{97}。
\fi
\end{checkpoint}

\subsection{默认路径的 \codet{past\_key\_values} 不是 cache，是 token 前缀}
\label{sec:b-kv-recompute}

Day 3--4 的 adapter 刻意\cnemph{不}做真 KV cache。\code{past_key_values} 里装的是到目前为止的
全部 token id，每一步 decode 都把这个前缀重新跑一遍完整前向。

\begin{examplebox}{adapter_call_trace}{同一个字段名，两种完全不同的东西}
把 \code{LMEngine} 驱动 adapter 的三步调用并排写出来。prompt 长度 5，生成 2 个 token。
左边是默认的 recompute 模式，右边是 \code{use_kv_cache=True}。

\begin{center}
\footnotesize
\setlength{\tabcolsep}{4pt}
\begin{tabular}{llll}
\toprule
调用 & 参数 & recompute 模式 & cache 模式 \\
\midrule
prefill
  & \codet{input\_ids} & \codet{(1,5)} 整个 prompt & \codet{(1,5)} 整个 prompt \\
  & \codet{past\_key\_values} & \codet{None} & \codet{None} \\
  & 模型内实际算 & 5 个位置 & 5 个位置 \\
  & 返回 \codet{.logits} & \codet{(1,5,V)} & \codet{(1,5,V)} \\
  & 返回 \codet{.past\_key\_values} & \codet{(1,5)} \cnemph{token id} & 16 层 $\times$ $(K,V)$，各 \codet{(1,12,5,64)} \\
\addlinespace
decode 1
  & \codet{input\_ids} & \codet{(1,1)} 新 token & \codet{(1,1)} 新 token \\
  & \codet{past\_key\_values} & \codet{(1,5)} & 上一步那 16 层 \\
  & 模型内实际算 & \cnemph{6 个位置}（全部重算） & \cnemph{1 个位置} \\
  & 返回 \codet{.logits} & \codet{(1,6,V)} & \codet{(1,1,V)} \\
  & 返回 \codet{.past\_key\_values} & \codet{(1,6)} token id & 16 层，各 \codet{(1,12,6,64)} \\
\addlinespace
decode 2
  & 模型内实际算 & \cnemph{7 个位置} & \cnemph{1 个位置} \\
  & 返回 \codet{.logits} & \codet{(1,7,V)} & \codet{(1,1,V)} \\
\bottomrule
\end{tabular}
\end{center}

\medskip
两件事一眼可见。第一，\code{.logits} 的形状在两种模式下\cnemph{不一样}：
recompute 返回整个前缀的 logits，cache 只返回新 token 的。这不影响 \code{LMEngine}，
因为它自己切 \code{[:, -1, :]}（\codeat{llm-serving/inferenceLM/engine/lm\_engine.py}{62}），
两种模式下这一切都取到"最后一个位置的 logits"。第二，累计到第 $n$ 步，
recompute 一共算了 $\sum_{i} (L{+}i) = O(n^2)$ 个位置，cache 是 $O(n)$。

形状是 16 层 $\times$ 12 头 $\times$ $d_k{=}64$，来自 \codeb{configs_sft/lr1e5_replay.json}
的 \codet{model} 段（\codet{num\_layers=16}, \codet{num\_heads=12}, \codet{d\_model=768}）。
\end{examplebox}

这个取舍是 \code{PLAN.md} 里就写好的，理由写进了模块 docstring。原话是：

\begin{lstlisting}[style=hlrepl]
Recompute keeps the two projects' seam small and honest: the serving layer is
genuinely running the locally-trained weights, and the cost of not having a cache
shows up as latency rather than as a correctness fudge.
\end{lstlisting}
\noindent 来源：\codeb{llm-serving/inferenceLM/engine/transformer_lm_adapter.py}:31--33，逐字手抄。

"体现为延迟而不是正确性上的糊弄"这句话是有内容的，不是修辞。它说的是：
recompute 解码和 cache 解码\cnemph{必须产生完全相同的 token 序列}，
差别只能出现在时间上。这条等式后来变成了 \secref{sec:b-kv-tests} 里所有测试的共同断言，
也变成了 \secref{sec:b-kv-flat} 那张对比表能被解释的前提——如果两条路径的输出会分叉，
比较它们的 tok/s 就没有意义。

\begin{checkpoint}{recompute_is_not_a_cache}{一个叫 cache 的字段里没有 cache}
不看代码回答：
\begin{enumerate}[label=(\alph*)]
  \item recompute 解码的复杂度是 $O(n^2)$，那 TTFT 会不会也受影响？
  \item 为什么 recompute 版本在 \codeb{prefix.shape[1] > n_positions} 时是从\cnemph{左边}截断？
  \item 这个设计在什么情况下会从"慢"变成"错"？
\end{enumerate}

\ifshowanswers
\tcblower
\answer (a) 不会。TTFT 是 prefill 那一次前向，两种模式下算的位置数完全相同
（\secref{sec:b-kv-recompute} 的例子第一行）。实测也印证了：
recompute 的 TTFT 是 61.3--65.7 ms，cache 是 59.3--63.0 ms，
差别落在 run-to-run 波动里（表~\ref{tab:kv-bench}）。KV cache 从定义上就不改善 TTFT。

(b) 因为 \code{LMEngine} 已经在 \code{inference()} 里拒绝了超过 \code{n_positions} 的
\code{max_length}，所以这条路径只有在有人\cnemph{绕过 \codet{LMEngine} 直接驱动 adapter}
时才会走到。这种情况下"丢掉最老的 token"是最不意外的行为——注释里写的是
"the least surprising behaviour"。

(c) 当 prefix 长到需要截断的时候。截断之后模型看到的位置从 0 重新开始编号，
生成会和"完整上下文"的结果不同——但这不是 cache 与 recompute 的差别，
是滑动窗口本身的语义。真 cache 路径在同样情况下的选择不一样：
\codeb{TransformerLM.forward_with_cache} 直接 \code{raise RuntimeError}，
理由是越界位置根本没有 RoPE 表项，静默回绕比报错糟糕得多
（\codeat{cs336-basics/cs336\_basics/lm.py}{203}）。\cnemph{两条路径在越界这件事上行为不一致}，
这是一个真实的不一致，讲义里记下来，没有修。

\whereincode \codeat{llm-serving/inferenceLM/engine/transformer\_lm\_adapter.py}{146}（左截断）、
\codeat{cs336-basics/cs336\_basics/lm.py}{203}（越界报错）。
\fi
\end{checkpoint}

\subsection{验证：把 GPT-2 自己剥成裸签名}
\label{sec:b-kv-stub}

recompute 解码"只会更慢、不会不同"是一个可以被证伪的断言。怎么证？

夜里用的办法是：拿 \code{GPT2LMHeadModel} 包一层，只暴露
\codeb{(batch, seq) -> (batch, seq, vocab)}——把它的 config、它的 cache、
它的输出 dataclass 全部剥掉，剥完之后它\cnemph{和我们自己的 \codet{TransformerLM} 长得一模一样}。
然后把这个裸模型塞进 adapter，用 \code{LMEngine} 跑贪心解码，
和 HuggingFace 自带的、\cnemph{带 KV cache 的} \code{generate()} 对拍。

\begin{lstlisting}[style=hlplain]
class _LogitsOnly(nn.Module):
    """Reduce a HuggingFace causal LM to the bare `(batch, seq) -> (batch, seq, vocab)`
    signature that a from-scratch `TransformerLM` has."""

    def __init__(self, hf_model: GPT2LMHeadModel) -> None:
        super().__init__()
        self.hf_model = hf_model

    def forward(self, input_ids: torch.Tensor) -> torch.Tensor:
        return self.hf_model(input_ids=input_ids).logits
\end{lstlisting}
\noindent 来源：\codeb{llm-serving/tests/transformer_lm_adapter_test.py}:15--30，手抄，
docstring 只保留了第一句。

结果是 token 级完全一致，6 个测试全过。这个手法的正式名字和它的适用边界在姊妹篇里：
\extref{a:differential-testing}。

\begin{checkpoint}{gpt2_as_stub}{为什么 stub 要用 GPT-2，而不是一个随机小模型}
不看代码回答：
\begin{enumerate}[label=(\alph*)]
  \item 用一个随机初始化的两层小模型当 stub，这个测试还成立吗？
  \item 这个测试到底隔离出了什么？它\cnemph{不能}证明什么？
\end{enumerate}
\hint 参照物是"HuggingFace 自己带 cache 的 \code{generate()}"。

\ifshowanswers
\tcblower
\answer (a) 部分成立但弱得多。随机小模型没有可信的参照实现——你只能拿它和它自己比，
那测的是"recompute 循环写得对不对"，不是"recompute 语义对不对"。
用 GPT-2 才能拿到一个\cnemph{独立实现的、带真 KV cache 的} \code{generate()} 当参照物：
两边如果对上，说明 adapter 的 recompute 语义和工业界公认的 cached 语义等价。
另外随机权重的 logits 很平，argmax 容易在数值噪声上翻转，会让测试变得又假阳又假阴。

(b) 隔离出来的是\cnemph{adapter 的 recompute 逻辑}——前缀怎么拼、\code{[:, -1, :]} 怎么切、
\code{past_key_values} 怎么传回去。因为模型两边都是同一个 GPT-2 权重，
模型本身的正确性被消掉了。它\cnemph{不能}证明的是我们自己的 \code{TransformerLM} 是对的
（那是训练那边的事），也不能证明真 KV cache 是对的——那是另一套测试，见
\secref{sec:b-kv-tests}。这两件事很容易混，但它们用的是完全不同的参照物：
这里的参照物是 HF 的 \code{generate()}，那边的参照物是我们自己的 \code{forward}。

\whereincode \codeat{llm-serving/tests/transformer\_lm\_adapter\_test.py}{62}。
\fi
\end{checkpoint}

\subsection{真 KV cache：三个新方法，\codet{forward} 一个字节没动}
\label{sec:b-kv-impl}

夜里第二轮才把真 cache 做出来，而且是刻意放在最后做的。实现方式是\cnemph{纯增量}：

\begin{itemize}[leftmargin=1.5em]
  \item \codeb{MultiHeadsAttention.forward_with_cache}（\codeat{cs336-basics/cs336\_basics/transfromer/multiheads\_attention.py}{127}）
  \item \codeb{PreNormTransformer.forward_with_cache}（\codeat{cs336-basics/cs336\_basics/transfromer/transformer.py}{77}）
  \item \codeb{TransformerLM.forward_with_cache}（\codeat{cs336-basics/cs336\_basics/lm.py}{183}）
\end{itemize}

三个\cnemph{新}方法。原来的 \code{forward} 一个字节没动，也就是\cnemph{训练路径零风险}——
改完之后在远端重跑 init check，输出和改之前逐位相同（W1 std 0.06533 / 初始 loss 7.0294 /
global grad norm 2.151，记录在 \codeb{interview_prep/PROGRESS.md}；本机没有对应日志可核对）。
把这条做成硬约束的代价是有一点重复代码（QKV 投影那几行在两个方法里各写了一遍），
换来的是"这个改动不可能影响到已经跑完的 7 小时训练"这个可以直接说出口的性质。

真 cache 为什么在数学上等价于重算——同一个位置的 $K$、$V$ 只依赖于它自己那个 token 的隐状态，
所以算过一次就可以存下来——这个论证属于姊妹篇，不在这里重复：\extref{a:kv-cache-legality}。
这一节只讲两个\cnemph{实现层面}的不变量，它们不是数学问题，是"写错了会静默出错"的问题。

\subsubsection{不变量 1：\codet{is\_causal} 不是手写开关，是从 shape 推出来的}

\lstinputlisting[style=hlnum, firstline=154, lastline=170, firstnumber=154]%
  {../../cs336-basics/cs336_basics/transfromer/multiheads_attention.py}

第 169 行是这一节的重点：

\begin{center}
\codeb{is_causal = Q_packed.shape[-2] == K_packed.shape[-2]}
\end{center}

没有 \code{if is_prefill:}，没有多传一个 \code{mode} 参数，也没有让调用方负责设对。
"什么时候该 mask"被表达成了一个\cnemph{关于 shape 的性质}：
query 和 key 数量相同 $\Leftrightarrow$ 它们是同一个 block $\Leftrightarrow$ 需要三角遮罩。
prefill 时 $\text{seq}_q = \text{seq}_k = n$，自动为真；decode 时 $\text{seq}_q = 1$、
$\text{seq}_k = n{+}1$，自动为假。

\begin{checkpoint}{is_causal_from_shape}{decode 时为什么必须 \codet{is\_causal=False}}
不看代码回答：
\begin{enumerate}[label=(\alph*)]
  \item 现有的 causal mask 是怎么构造的？单 query 时它会露出哪些 key？
  \item 如果 decode 时错误地传了 \code{is_causal=True}，模型会崩溃、还是会输出看起来正常但错误的结果？
  \item 用 shape 推导代替一个显式开关，除了少写一个参数，还买到了什么？
\end{enumerate}
\hint mask 的形状是 $(\text{seq}_q, \text{seq}_k)$，不是 $(\text{seq}_k, \text{seq}_k)$。

\ifshowanswers
\tcblower
\answer (a) \codeb{causal_mask = torch.tril(torch.ones(seq_q, seq_k))}
（\codeat{cs336-basics/cs336\_basics/transfromer/scaled\_dot\_prod\_attention.py}{49}）。
这是一个\cnemph{左对齐}的下三角：第 $i$ 行露出 key $0 \ldots i$。
decode 时 $\text{seq}_q = 1$，只有第 0 行，于是\cnemph{只露出 key 0}——
整个缓存里除了第一个 token 全被 mask 掉。

(b) \cnemph{不会崩溃}。形状全部合法，softmax 正常，输出的 logits 形状也对。
模型会安静地生成一串"只看得见第一个 token"的续写。这正是这一类 bug 最危险的地方：
它不产生任何异常，只产生更差的文本，而一个 190M 的欠训模型本来输出就不怎么样，
很容易被解释成"模型就是这个水平"。

(c) 买到的是\cnemph{这个 bug 不可能通过传错参数发生}。显式开关意味着每一个调用点
（三层 \code{forward_with_cache} 逐层往下传）都是一次犯错的机会，而且传错了没有任何检查。
shape 推导把它降级成一个局部的、不依赖调用方的判断。代价是它对"seq\_q 恰好等于 seq\_k
但语义上不是同一个 block"这种情形会判错——本项目里不存在这种调用，
但这是这个技巧的真实边界。

\whereincode \codeat{cs336-basics/cs336\_basics/transfromer/multiheads\_attention.py}{169}（推导），
\codeat{cs336-basics/cs336\_basics/transfromer/multiheads\_attention.py}{165} 起的三行注释写着为什么。
mask 构造在 \codeat{cs336-basics/cs336\_basics/transfromer/scaled\_dot\_prod\_attention.py}{46}--50。
\fi
\end{checkpoint}

\subsubsection{不变量 2：RoPE 必须按绝对位置旋转后再进 cache}

上面那段代码的第 154--155 行注释说的就是这件事：
"Append to the cache AFTER rotating: RoPE is applied once, at each token's own absolute
position, so cached keys must already carry their rotation."

绝对位置是在最外层算的：

\lstinputlisting[style=hlnum, firstline=201, lastline=213, firstnumber=201]%
  {../../cs336-basics/cs336_basics/lm.py}

\code{past_len} 从缓存的形状里读出来（第 201 行），positions 从 \code{past_len} 开始切
（第 212 行）。所以第 $n$ 步 decode 的那个 token 拿到的是位置 $n$，不是位置 0。

\begin{checkpoint}{rope_before_cache}{把旋转和入缓存的顺序颠倒，哪个测试会抓到}
不看代码回答：
\begin{enumerate}[label=(\alph*)]
  \item 如果 \code{positions} 写成 \codeb{self.token_positions[:seq_new]}（chunk 内相对位置），
        \code{tests/test_kv_cache.py} 的 6 个测试里哪些会红？
  \item 哪些\cnemph{不会}红？为什么？
  \item 这个 bug 在 prefill 阶段会不会暴露？
\end{enumerate}
\hint prefill 时 \code{past_len} 是 0。

\ifshowanswers
\tcblower
\answer 这三问是写讲义时\cnemph{真的注入这个 bug 跑了一遍}得到的，不是推理
（做法见表~\ref{tab:kv-mutation} 的 caption）。

(a) 两个会红：\codeb{test_greedy_decoding_is_token_identical} 和
\codeb{test_incremental_logits_match_full_forward_at_each_step}。

(b) 四个不会红。\codeb{test_prefill_matches_plain_forward} 不红是因为 prefill 时
\code{past_len=0}，\code{[:seq_new]} 和 \code{[0:seq_new]} 是同一个切片——
\cnemph{bug 在 prefill 路径上根本不激活}。
\codeb{test_cache_grows_by_one_key_per_step} 和
\codeb{test_rejects_decoding_past_the_context_length} 只断言形状和边界，位置旋转对它们不可见。
\codeb{test_batched_decoding_matches_per_sequence_decoding} 不红最值得注意：
批量和逐条\cnemph{犯的是同一个错}，两边一起错、错得一样，所以它们仍然相等。
\cnemph{一个测试只能抓到破坏它所断言的那个对称性的 bug}。

(c) 不会。这就是 (b) 的第一句。这个 bug 只在"缓存非空"时激活，
所以任何只测 prefill 的检查都是安全通过的——包括
\code{kv_cache_gate.py} 的第一项 \code{check_logits_match}。
门禁能挡住它，靠的是第二项 \code{check_greedy_match}。

\whereincode 位置切片在 \codeat{cs336-basics/cs336\_basics/lm.py}{212}；
旋转-再入缓存的顺序在 \codeat{cs336-basics/cs336\_basics/transfromer/multiheads\_attention.py}{150}--159。
\fi
\end{checkpoint}

\subsection{怎么测一个 cache}
\label{sec:b-kv-tests}

\code{tests/test_kv_cache.py} 有 6 个测试。它们不是六件事，是\cnemph{同一件事的六个角度}：
cache 是一个优化，所以唯一可接受的结果是\cnemph{它只改变速度}。参照物是 \code{forward}，
理由写在文件的 docstring 里：那是模型被训练时跑的那条代码，是唯一可信的基准。

写这份讲义时，我把这 6 个测试拿去做了一次\cnemph{变异测试}：往 cache 实现里逐个注入
五类真实的 bug，看哪些测试会红。结果如下。

\begin{table}[htbp]
\centering
\small
\setlength{\tabcolsep}{5pt}
\begin{tabular}{lccccc}
\toprule
注入的 bug & T1 & T2 & T3 & T4 & T5/T6 \\
\midrule
RoPE 按 chunk 内相对位置旋转 & --- & 抓到 & 抓到 & --- & ---/--- \\
\codet{is\_causal} 硬编码 \codet{True} & --- & 抓到 & 抓到 & --- & ---/--- \\
\codet{is\_causal} 硬编码 \codet{False} & 抓到 & 抓到 & 抓到 & --- & ---/--- \\
缓存 key 跨 batch 行串用 & --- & --- & --- & --- & ---/抓到 \\
缓存 key 有 0.1\% 数值偏差 & --- & --- & 抓到 & --- & ---/--- \\
\bottomrule
\end{tabular}
\caption{五类注入 bug $\times$ 6 个测试。
T1 $=$ \codeb{test_prefill_matches_plain_forward}，
T2 $=$ \codeb{test_greedy_decoding_is_token_identical}，
T3 $=$ \codeb{test_incremental_logits_match_full_forward_at_each_step}，
T4 $=$ \codeb{test_cache_grows_by_one_key_per_step}，
T5 $=$ \codeb{test_rejects_decoding_past_the_context_length}，
T6 $=$ \codeb{test_batched_decoding_matches_per_sequence_decoding}。
出处：\cnemph{写这份讲义时现做的实验，夜里没有做过}——用 monkeypatch 替换
\codeb{MultiHeadsAttention.forward_with_cache} 和 \codeb{TransformerLM.forward_with_cache}
注入每一类 bug，再对\cnemph{未经修改的} \codet{tests/test\_kv\_cache.py} 跑 pytest。
本机 CPU，每轮约 0.15 秒。未注入时 6 passed。}
\label{tab:kv-mutation}
\end{table}

从这张表能读出三件事，其中两件是设计时想到的，一件不是。

\cnemph{T4 和 T5 没有抓到任何一个注入的 bug。} 这不是它们没用——它们守的是形状和边界，
而我注入的五类 bug 全是数值/语义的。一个把缓存拼错方向（\code{dim=-1} 而不是 \code{dim=-2}）
的 bug 会立刻被 T4 抓到。但这张表确实说明：\cnemph{这套测试的实际防线是 T3 和 T6}。

\cnemph{T6 是唯一能抓到批内串 key 的测试。} 其他五个测试全部用 batch=1，
在那个维度上没有任何断言。

\cnemph{T3 比 T2 严格。} 最后一行是专门为验证这一点造的：给缓存里的 key 乘 1.001，
logits 会动，但动得不足以翻转任何一个 argmax——于是 T2（比 token）全过，T3（比 logits）红。

\begin{checkpoint}{stepwise_comparison}{为什么要逐步比对，而不是只看最终状态}
不看代码回答：
\begin{enumerate}[label=(\alph*)]
  \item T3 每一步都拿 \code{model(ids)} 重跑一遍完整前向来对照。这很慢，为什么值得？
  \item 表~\ref{tab:kv-mutation} 最后一行说明 T3 比 T2 敏感在哪？
  \item 如果只留一个测试，留哪个？
\end{enumerate}
\hint T3 的 docstring 自己给了答案的一半。

\ifshowanswers
\tcblower
\answer (a) 因为"陈旧的 key"和"错位置的旋转"这两类 bug 的共同特点是\cnemph{要解码若干步之后才显形}。
只检查最终状态、而且是在一个短序列上检查，很可能整段都还没进入 bug 激活的区间。
T3 的 docstring 原话是 "a cache bug that only shows up deep into decoding
(a stale key, a mis-rotated position) would pass an end-state check on a short run"。
慢是真的慢（10 步就要跑 11 次完整前向），但这是一个 3 层 $d_{\text{model}}{=}64$ 的
CPU 模型，整份文件跑完不到 1 秒。

(b) 敏感在\cnemph{比较的对象是连续量而不是离散量}。T2 比的是 argmax 之后的 token id，
一个小到不足以翻转 argmax 的数值误差它看不见；T3 比的是 logits 本身，
\code{atol=1e-4}。注入 0.1\% 的 key 偏差时 T2 全过、T3 红，就是这个差别的直接证据。
反过来说 T2 也不是多余的：它断言的是\cnemph{用户真正看得到的那个性质}——
两条路径生成的文本必须一模一样。T3 red 而 T2 green 时，你知道有数值漂移但输出还没变；
两个都 red 时，问题已经跑到用户面前了。

(c) 留 T3。它是唯一一个既覆盖 prefill 又覆盖多步 decode、又比较连续量的测试。
但留了 T3 之后就没有任何测试守着 batch 维度（表~\ref{tab:kv-mutation} 第四行），
所以实际上不能只留一个。

\whereincode \codeat{tests/test\_kv\_cache.py}{66}。
\fi
\end{checkpoint}

\lstinputlisting[style=hlnum, firstline=66, lastline=77, firstnumber=66]%
  {../../tests/test_kv_cache.py}

\begin{checkpoint}{batched_decode_leak}{批内串 key 是什么样的 bug}
不看代码回答：
\begin{enumerate}[label=(\alph*)]
  \item T6 拿 batch=2 的解码结果和逐条解码的结果比。这能抓到什么？
  \item 为什么其他五个测试抓不到？
  \item 本项目的 \code{LMEngine} 目前一次只处理一个请求，为什么还要写这个测试？
\end{enumerate}

\ifshowanswers
\tcblower
\answer (a) 抓到任何"缓存记账不是逐序列的"错误：把一行的 $K$ 广播给整批、
reshape 时把 batch 和 head 两个维度搞混、cache 拼接用错了 dim。
这一类 bug 的表现是 batch 里第 $i$ 条请求看到了第 $j$ 条请求的历史——
在真实 serving 里这既是错误也是数据泄漏。变异实验里我注入的是最直白的一种
（\codeb{past_k[0:1].expand_as(past_k)}），只有 T6 变红。

(b) 因为其他五个都是 batch=1。batch=1 时"跨行串用"和"不串用"是同一件事，
没有可以被破坏的对称性。这是表~\ref{tab:kv-mutation} 想说的核心：
\cnemph{测试能抓到的，只有它所断言的那个不变量被破坏的情形}，
再多的测试也不会自动覆盖一个没人断言过的维度。

(c) 因为这是给\cnemph{还没写的那个调度器}提前立的规矩。
\code{llm-serving} 的路线图上 continuous batching 是明确的下一步
（见姊妹篇 \extref{a:continuous-batching-def} 和 \extref{a:stage-homogeneous}）。
到那时候 batch 维度会突然变成热路径，而那时再补这个测试，
就得先在一个已经变复杂的系统里定位 bug。现在写，成本是十几行。

\whereincode \codeat{tests/test\_kv\_cache.py}{104}。
\fi
\end{checkpoint}

\begin{tipbox}{不用 GPU 也能验 cache}
\code{tests/test_kv_cache.py} 整份文件用的是 CPU 上一个
\codeb{vocab_size=257, context_length=64, num_layers=3, d_model=64} 的模型。
本机实测 \code{6 passed in 0.96s}（首次运行含 import；变异实验里热跑是 0.15 秒）。

这不是"降配将就"，这是这类测试的\cnemph{正确规模}：
\cnemph{cache 的正确性和模型规模无关}。它断言的是一个代数等式——
带缓存的前向必须和不带缓存的前向给出同样的数——而这个等式在
$d_{\text{model}}{=}64$ 和 $d_{\text{model}}{=}768$ 上是同一个等式。
规模只影响\cnemph{性能}结论（\secref{sec:b-kv-flat} 那张表就必须在真 GPU、真权重上测），
不影响正确性结论。

顺带说，这条性质是免费拿到的吗？不是。\code{cs336_basics} 里有 8 个文件
无条件 \codeb{import torch.cuda.nvtx as nvtx} 并用 \code{@nvtx.range(...)} 标注 forward，
在非 CUDA build 上导入正常、一调用就 \code{RuntimeError}。
是先加了 \codeb{cs336_basics/nvtx_compat.py}（import 时探测一次，不可用就退化成 no-op），
\code{TransformerLM} 才能在这台笔记本上跑起来。
\resources CPU 即可；\codeb{uv run pytest tests/test_kv_cache.py}，约 1 秒。
\end{tipbox}

除了进常规测试套件的这 6 个测试，还有一个一次性的门禁脚本
\codeb{cs336_systems/experiments/kv_cache_gate.py}：缓存解码必须和重算解码
\cnemph{logits 逐位相同 + 贪心 token 完全一致}，否则 \code{sys.exit(1)}，
调用方不许启用 cache。

\begin{examplebox}{kv_gate_output}{门禁的真实输出}
夜间窗口 stage 7 的日志原文（\codeb{run_logs/post_train_window.log}:1123--1132）：

\begin{lstlisting}[style=hlshell]
[kv] gating against checkpoints_pretrain_v1/checkpoint_epoch_0100.pt
=== gate 1: small random model ===
  prefill logits max|delta| = 0.000e+00  -> PASS
  greedy tokens identical    -> PASS

=== gate 2: real trained model ===
  prefill logits max|delta| = 0.000e+00  -> PASS
  greedy tokens identical    -> PASS

GATE PASSED: cached decode is token-identical to recompute.
\end{lstlisting}

\code{0.000e+00} 是逐位相同，不是"在容差内"。这比预期更强：
门禁自己的容差是 \code{tol=2e-4}，脚本 docstring 里也写着"不同的 reduction order
产生微小 logit 差异是预期的"。实际拿到 0 是因为 prefill 路径上
\code{forward_with_cache} 和 \code{forward} 做的是\cnemph{同样顺序的同样运算}
（\code{past_kv=None} 时那两个分支的算子序列一致），所以浮点结果逐位一致。
真正需要容差的是 decode 步——但那里比的是 token 不是 logits。
\end{examplebox}

\begin{checkpoint}{gate_with_final_weights}{为什么门禁要拿最终权重再跑一次}
不看代码回答：
\begin{enumerate}[label=(\alph*)]
  \item 小随机模型已经过了，为什么还要拿 190M 的真 checkpoint 再过一次？
  \item 门禁脚本和 \code{tests/test_kv_cache.py} 有什么分工？为什么不合并？
\end{enumerate}

\ifshowanswers
\tcblower
\answer (a) 因为\cnemph{拿别的权重过的门禁不算这批权重的门禁}。原则上 cache 的正确性
和权重无关（\secref{sec:b-kv-tests} 的 tipbox 就是这么论证的），但"原则上无关"
和"这批要被报告数字的权重上确实无关"是两个命题，后者只要 30 秒就能验，没有理由不验。
更实际的理由是：真 checkpoint 会额外走一遍 \codeb{k.removeprefix("_orig_mod.")} 的加载路径、
真实的 \code{context_length=1024} 的 RoPE 表、真实的 16 层深度——
这些都是小模型上不存在的东西。门禁通过之后才允许 benchmark 报缓存的数字，
顺序在夜间脚本里是硬的：gate 失败则 \code{sys.exit(1)}，后面的 \code{bench_serving.py} 不会跑。

(b) 分工是"常驻回归"与"一次性放行"。\code{test_kv_cache.py} 进常规套件，
每次 \code{uv run pytest} 都跑，用小模型、跑得起，防的是\cnemph{将来的改动}打破 cache。
\code{kv_cache_gate.py} 是一个带 \code{--checkpoint} 参数的脚本，用真权重、
在 GPU 机器上跑，防的是\cnemph{这一次要报的数字}建立在没验过的路径上。
合并的话要么让常规套件依赖一个 2.3GB 的 checkpoint（跑不起来），
要么让门禁失去"针对特定权重"的语义。

\whereincode \codeat{cs336\_systems/experiments/kv\_cache\_gate.py}{94}（gate 2 分支）、
\codeat{cs336\_systems/experiments/kv\_cache\_gate.py}{118}（失败即退出）。
\fi
\end{checkpoint}

\subsection{KV cache 的收益和预期完全不一样}
\label{sec:b-kv-flat}

门禁过了，可以报数字了。

\begin{table}[htbp]
\centering
\small
\setlength{\tabcolsep}{5pt}
\begin{tabular}{rrrrrrrr}
\toprule
& \multicolumn{3}{c}{recompute} & \multicolumn{3}{c}{cache} & \\
\cmidrule(lr){2-4}\cmidrule(lr){5-7}
prompt tokens & TTFT ms & p50 ms & tok/s & TTFT ms & p50 ms & tok/s & 加速比 \\
\midrule
16  & 61.3 & 63.03 & 15.9 & 62.3 & 52.45 & 19.2 & $1.20\times$ \\
64  & 65.7 & 61.25 & 16.4 & 63.0 & 48.06 & 21.0 & $1.28\times$ \\
128 & 61.8 & 62.79 & 15.7 & 59.3 & 52.78 & 18.9 & $1.20\times$ \\
256 & 65.3 & 64.08 & 15.6 & 61.2 & 51.80 & 19.8 & $1.27\times$ \\
512 & 64.5 & 64.82 & 15.6 & 61.5 & 52.12 & 19.2 & $1.23\times$ \\
\bottomrule
\end{tabular}
\caption{recompute 与真 KV cache 的完整对照。
出处：\codeb{artifacts/serving_benchmark_kv.json}，
单卡 \codet{cuda:0}、0.19B 模型、fp32、batch 1、\codet{new\_tokens=64}（即每格 63 个 decode 步）、
两种模式用同一个 \codeb{torch.Generator().manual_seed(0)} 生成\cnemph{逐字节相同}的 prompt、
2 次 warmup。p50 是 63 个 step 时间的中位数。加速比 $=$ cache tok/s $\div$ recompute tok/s，
取自 JSON 的 \codeb{cache_speedup_by_prompt_len} 字段。
LMEngine 端到端往返（22 token）：recompute 1.36 s，cache 1.14 s。}
\label{tab:kv-bench}
\end{table}

五个加速比：$1.20 / 1.28 / 1.20 / 1.27 / 1.23$。均值 $1.24\times$。
\cnemph{prompt 长度涨了 32 倍，加速比几乎不动}（末项比首项 $=1.02$）。

这和事前的断言完全相反。事前的断言写在下面这个框里，原样保留：

\begin{retraction}{加速比会随 prompt 长度增长，因为重算是 \codet{O(prefix)} 而 cache 是 \codet{O(1)}}
这句话的\cnemph{前半段推理是对的}——重算解码每步确实要跑完整个前缀，
cache 每步确实只跑一个位置。错的是从这个渐进关系直接跳到"所以在这台机器、这个模型上，
加速比会随长度增长"。渐进复杂度描述的是当 prefix 大到能主导总成本时的行为；
它不告诉你在 prefix $=512$、模型 190M、batch $=1$ 的时候，
prefix 相关的那部分占总成本的多少。

实测占比很小，所以加速比是平的。这条被 \codeb{artifacts/serving_benchmark_kv.json}
的五个数直接否掉，撤回。
\end{retraction}

原因不难说清楚，但要说得准确需要一点算术。recompute 模式下，
prompt 从 16 涨到 512（32 倍）时，每步 p50 从 63.03 ms 涨到 64.82 ms——
\cnemph{只涨了 2.8\%}。也就是说\cnemph{连重算本身都几乎不随前缀长度增长}。
在这个规模上，每一步的成本几乎全部是"把 16 层权重从头到尾过一遍"这个\cnemph{固定开销}，
而 KV cache 消掉的那部分（前缀相关的 attention 计算和它的 QKV 投影）只是其中一小块。

这个固定开销到底是什么，值得记一笔\cnemph{没有查证到底}的地方。夜报里的措辞是
"dominated by the fixed forward pass over the weights"，暗示是访存。粗算一下：
0.19B 参数 $\times$ 4 字节 $= 760$ MB，RTX 3090 的显存带宽约 936 GB/s，
把权重读一遍的下界是 \cnemph{约 0.8 ms}；而实测每步 60 ms，
是这个下界的 \cnemph{约 70 倍}。计算量更不是瓶颈（batch 1 单 token 约 0.38 GFLOP，
在 fp32 的 3090 上是 0.01 ms 量级）。所以这 60 ms 里绝大部分\cnemph{既不是访存也不是算力}，
更像是每层十几次小 kernel launch、einops 的 \code{rearrange}、Python 层调度的累加。
\cnemph{这只是一个数量级估算，没有做过 profiling}，所以只能作为一个待查的开放问题写在这里。
但它不影响本节的结论：不论那 60 ms 具体花在哪，它都是\cnemph{与前缀长度无关}的固定成本，
而这正是加速比平坦的直接原因。

还有一件必须说的：这张表跑的是\cnemph{朴素的 PyTorch attention，不是本项目的 Triton kernel}。
\codeb{bench_serving.build_adapted_model} 构造 \code{TransformerLM} 时没有传
\code{attention_fn}，走的是默认值 \codeb{scaled_dot_product_attention}
（\codeat{cs336-basics/cs336\_basics/lm.py}{131}）；
而 \code{forward_with_cache} 那条路径更是\cnemph{直接写死}调用
\codeb{scaled_dot_product_attention}（\codeat{cs336-basics/cs336\_basics/transfromer/multiheads\_attention.py}{170}），
连 \code{self.attention_fn} 都不看。两种模式因此是同一个 attention 实现，
比较是公平的；但\cnemph{这张表和 Triton kernel 无关}，不能拿它去说 kernel 的任何事。

\begin{checkpoint}{flat_speedup}{加速比为什么是平的}
不看代码回答：
\begin{enumerate}[label=(\alph*)]
  \item 在什么条件下 KV cache 才会给出教科书上那种数量级的加速？
  \item 表~\ref{tab:kv-bench} 里哪一列最能说明"重算不贵"？
  \item $1.24\times$ 这个数字，能不能写进简历说"实现了 KV cache，decode 提速 24\%"？
\end{enumerate}

\ifshowanswers
\tcblower
\answer (a) 当前缀相关的工作在每步成本里占比变大的时候，也就是：上下文更长
（attention 是 $O(n)$ 每步，权重那一遍是 $O(1)$，$n$ 大到一定程度必然反超）、
batch 更大（权重那一遍在整个 batch 上摊薄，而 KV 的量随 batch 线性增长）、
或者模型更小。这三条在这次的配置里全是反的：$n \le 576$、batch $=1$、190M 参数 fp32。
\cnemph{KV cache 在这个规模上不划算}——这是个比"cache 快 N 倍"更有信息量的结论。

(b) recompute 那一列的 p50：$63.03 \to 61.25 \to 62.79 \to 64.08 \to 64.82$ ms。
前缀涨了 32 倍，每步只贵了 2.8\%。如果重算真的贵，这一列应该是明显上翘的。
（JSON 里 \codeb{decode_cost_growth.latency_ratio} 字段记的就是这个比值：$1.028$。）

(c) 可以，但必须带上后半句："在 190M/fp32/batch=1/上下文 $\le$ 576 的条件下测得
$1.24\times$，而且加速比不随 prompt 长度增长，因为这个规模下每步成本由固定的权重前向主导。"
只报 $1.24\times$ 会让人以为这是个一般结论；只报"实现了 KV cache"则回避了最有意思的部分。
真正值得讲的是"我预期它随长度增长，实测是平的，我去查了为什么"这个过程。
\fi
\end{checkpoint}

\subsubsection{一个在读自己数据之前就写好结论的 benchmark}

上面那句被撤回的话，不是我在报告里随手写的一句话。它是
\code{bench_serving.py} 里一个\cnemph{无条件 \codet{print}} 出来的硬编码结论，
不看任何实际数字。\code{make_night_report.py} 里也有同一句的 markdown 版本。

后果在夜间日志里留着，原样引在这里
（\codeb{run_logs/post_train_window.log}:1167--1175）：

\begin{lstlisting}[style=hlshell]
=== recompute vs KV cache ===
  prompt   recompute tok/s   cache tok/s   speedup
      16              15.9          19.2     1.20x
      64              16.4          21.0     1.28x
     128              15.7          18.9     1.20x
     256              15.6          19.8     1.27x
     512              15.6          19.2     1.23x

The speedup grows with prompt length because recompute is O(prefix) per step and the
cache is O(1) in the prefix -- that widening gap IS the cost of not having a cache.
\end{lstlisting}

一张明显没有趋势的表，紧接着一行断言它在增长。两处都改成了按实测算趋势再措辞：
$>1.25\times$ 说"增长"，$<0.8\times$ 说"反常，需要查"，否则说"基本持平在 $\sim X\times$"。
修复后的代码和它的注释在本仓库里，可以直接引：

\lstinputlisting[style=hlnum, firstline=181, lastline=205, firstnumber=181]%
  {../../cs336_systems/experiments/bench_serving.py}

第 183--184 行那句 "A benchmark that states its conclusion before reading its own numbers
is worse than no benchmark" 是这一整章里最该记住的一句。

\begin{checkpoint}{benchmark_wrote_conclusion_first}{为什么这类 bug 比测错数字更严重}
不看代码回答：
\begin{enumerate}[label=(\alph*)]
  \item 一个 benchmark 把数字测错了，和一个 benchmark 数字测对了但结论是硬编码的，
        哪个更危险？
  \item 修复为什么要写成三个分支（增长 / 反常 / 持平），而不是简单地把那句话删掉？
  \item 这个 bug 是怎么被发现的？如果没被发现，会怎样？
\end{enumerate}

\ifshowanswers
\tcblower
\answer (a) 后者。测错数字的 benchmark 迟早会和别的证据对不上，而且它的失败是\cnemph{随机方向}的。
硬编码结论的 benchmark 每一次运行都朝\cnemph{同一个方向}撒谎，而且它撒的谎恰好是你最想听的
那一句——因为那句话正是你写它的时候相信的。它把 benchmark 从"检验假设的工具"
降级成"把假设重复一遍的机器"。更糟的是它\cnemph{看起来在工作}：表格是真的，
数字是真的，只有那一行结论是假的。

(b) 因为删掉之后，读日志的人要自己从五个数里看出趋势，而这正是当初想省掉的那一步。
三个分支保留了"脚本替你读表"这个便利，同时把"读的结果"和"数据"绑定起来。
更重要的是 \code{<0.8×} 那个分支：它守的是一个\cnemph{我不希望看到但必须能看到}的情形
（cache 比重算还慢，说明实现有问题），措辞是"worth investigating before quoting these
numbers"。一个只会说好消息的脚本，和一个硬编码结论的脚本，性质是一样的。

(c) 是在把表格和结论并排读的时候发现的——五个数明显没趋势，下一行说它在增长，
荒唐到不用算。如果没发现，\code{NIGHT_REPORT.md} 里会留下一句和自己表格矛盾的结论，
而这份报告是给别人看的。同一个晚上还在 bucket size 扫描的汇总代码里犯了同一类错
（直接 \code{print} "best: 50MB"，见 \secref{sec:b-bench}），说明这不是一次手滑，
是一种习惯：\cnemph{把"我预期会看到什么"写进了本该报告"我看到了什么"的地方}。

\whereincode \codeat{cs336\_systems/experiments/bench\_serving.py}{181}（注释与三分支）、
\codeat{cs336\_systems/experiments/make\_night\_report.py}{166}（同一处修复的 markdown 版本）。
\fi
\end{checkpoint}

\begin{misconception}{\codeb{artifacts/serving_benchmark_kv.json} 是修复后的脚本产出的，所以它是干净的。}
不是。那个 JSON 写于 2026-07-27 23:58（日志时间戳），修复是之后做的。
证据：修复后的代码会往结果里写一个 \codeb{cache_speedup_trend_last_over_first} 字段
（\codeat{cs336\_systems/experiments/bench\_serving.py}{188}），
而磁盘上那份 JSON 的顶层 key 只有
\codeb{device / new_tokens / modes / cache_speedup_by_prompt_len}，\cnemph{没有那个字段}。

这不影响任何数字——五个加速比是同一次测量的产物，修复只动了措辞逻辑。
但它意味着\cnemph{那份 artifact 是用有 bug 的脚本产出的}，
\code{NIGHT_REPORT.md} 里正确的"roughly flat at ~1.24x"那句话来自
\code{make_night_report.py} 事后重新生成。这一条写在这里，是因为
"结论对了"和"产出结论的代码对了"是两件事，而这一章刚好在讲这个区别。
\end{misconception}

\subsection{端到端：链路里除 tokenizer 外没有 HuggingFace 模型}
\label{sec:b-kv-e2e}

\codeb{tests/trained_model_integration_test.py} 走完整条链路：
\code{RequestReceiver} $\to$ \code{pending_queue} $\to$ \code{InferenceEngine} $\to$
\code{LMEngine} $\to$ \code{TransformerLMAdapter} $\to$ 自训权重。
2026-07-27 夜里在 GPU 机器上真跑通了，2 passed。输出：

\begin{lstlisting}[style=hlshell]
prompt:    The history of the internet began in the
continued:  late 19th century, and the first wave of internet-connected websites
           emerged in the 20th century.
\end{lstlisting}

语法通顺、切题，事实是错的（还输出过 "Paris is the capital \ldots\ of the United States"），
并且开始复读。一个 190M 参数、只训了 2.8 token/param 的模型就该是这个样子，
\cnemph{不要把它说得比实际好}。这段输出记录在 \codeb{interview_prep/PROGRESS.md}；
本机没有对应的测试日志，属于核不到出处的一项。

链路里除了 tokenizer，没有任何 HuggingFace 模型。而 tokenizer 之所以能直接用 HF 的，
是因为训练时用的 \codeb{tiktoken.get_encoding("gpt2")} 和 HF 的 \code{gpt2} tokenizer
是 token-for-token 相同的。

\begin{checkpoint}{tokenizer_compat}{凭什么说两个 tokenizer 是同一个}
不看代码回答：
\begin{enumerate}[label=(\alph*)]
  \item 训练用 tiktoken、serving 用 \code{AutoTokenizer}，怎么保证 token id 对得上？
  \item 如果不对得上，会表现成什么样？会不会报错？
  \item 项目里为什么最后没有自己训 BPE？
\end{enumerate}

\ifshowanswers
\tcblower
\answer (a) 两者是同一份 GPT-2 BPE：50257 个词表项，\codet{<|endoftext|>} 都是 50256。
本机核对过 tiktoken 这一半：\codeb{tiktoken.get_encoding("gpt2")} 的
\code{n_vocab=50257}、\code{eot_token=50256}，和 adapter 里那个
\code{GPT2_EOS_TOKEN_ID = 50256} 常量一致
（\codeat{llm-serving/inferenceLM/engine/transformer\_lm\_adapter.py}{47}）。
\cnemph{HF 那一半本机没有 \codet{transformers} 可以核对}，
所以"token-for-token 相同"这个说法目前的依据是 adapter docstring 第 35--38 行
和集成测试第 75 行的注释，加上两边都声称是 GPT-2 的原始 BPE。
真正把它验成事实的是集成测试本身：如果 id 不一致，
用 HF tokenizer 编码的 prompt 喂进按 tiktoken 训的模型，输出会是乱码而不是通顺英文。

(b) \cnemph{不会报错}，因为两边 vocab size 相同，任何 id 都在合法范围内。
表现是模型输出彻底的胡言乱语——但一个欠训的 190M 模型输出本来就不太好，
这两种"不好"很容易混淆。这正是为什么集成测试要 \code{print} 出续写文本让人眼看，
而不只是断言 \code{len(generated) > 0}。

(c) 完整的理由属于预训练那一章（\secref{subsec:tokenizer}）：自训 BPE 在真实混合语料上
慢到会把整个时间预算吃光，于是换成 tiktoken 的预训练 \code{gpt2} 编码，
780MB + 20MB 语料 92 秒编码完（\codeb{run_logs/tiktoken_encode.log}：86.0 s + 5.9 s）。
这里只记它\cnemph{顺带的收益}：serving 端因此不需要分发一份自定义 tokenizer，
\code{RequestReceiver} 的 \code{AutoTokenizer} 直接指向 \code{"gpt2"} 就行，一行都不用改。

\whereincode \codeat{llm-serving/tests/trained\_model\_integration\_test.py}{75}。
\fi
\end{checkpoint}

\subsection{环境这一坑}
\label{sec:b-kv-env}

最后记一件纯工程的事，因为它花掉的时间不比上面任何一节少。

要在 GPU 机器上同时跑两个仓库的代码，第一反应是给 \code{llm-serving} 单独
\code{uv sync} 一份 venv。查了之后发现\cnemph{两边 venv 各缺一半}：

\begin{itemize}[leftmargin=1.5em]
  \item \code{llm-serving} 没有 \code{einops} 和 \code{tiktoken}，
        而这两个是 \code{cs336_basics} 的必需依赖；
  \item 本项目没有 \code{transformers} 和 \code{pytest-asyncio}，
        而 \code{llm-serving} 的测试全靠它们。
\end{itemize}

而且 \code{llm-serving} 的 \code{pyproject.toml} pin 的是 \code{torch>=2.11}，
单独装会再拉一份约 3GB 的 torch（本项目跑的是已经验证过的 2.6），
而那台实例的根分区只剩约 10GB。

最后的解法是\cnemph{不建第二个 venv}：复用本项目的 venv，
用 \codeb{uv run --with transformers --with pytest-asyncio} 临时补上缺的两个包，
再用 \code{PYTHONPATH=~/llm-serving} 让 \code{inferenceLM} 可导入。
不改任何 \code{pyproject.toml}，torch 版本不变。

\begin{lstlisting}[style=hlshell]
PYTHONPATH=~/llm-serving uv run --with transformers --with pytest-asyncio \
  python cs336_systems/experiments/bench_serving.py \
    --config cs336_systems/experiments/real_pretrain_config.json \
    --checkpoint checkpoints_pretrain_v1/checkpoint_epoch_0100.pt \
    --modes recompute cache
\end{lstlisting}

这里有一个反直觉的细节，踩过一次才知道：\cnemph{\codet{uv pip install} 装的包活不过下一次
\codet{uv run}}。\code{uv run} 会把 venv 重新同步回 lockfile，把手动装的包清掉。
所以必须用 \code{--with}（它每次都重新解析这几个额外依赖），
而不是先 \code{uv pip install} 再 \code{uv run}。

这一坑之所以值得写进讲义，是因为它是"两个仓库解耦"这个设计决定的\cnemph{账单}。
\secref{sec:b-kv-adapter} 说 adapter 不 import \code{cs336_basics} 换来了三个好处；
代价就在这里——总有一个地方要把两边拼起来，那个地方是运行环境，
而运行环境比 import 语句难 debug 得多。这个账单是划算的，
但把它记下来比假装它不存在诚实。
